\subsection{Kết luận chương 2}

Chương 2 đã hoàn thành ba nhiệm vụ cốt lõi của giai đoạn đặc tả, phân tích và thiết kế hệ thống bỏ phiếu điện tử: làm rõ bài toán cần giải quyết, mô hình hóa các yêu cầu thành các ca sử dụng cụ thể, và chuyển hóa các ca sử dụng đó thành hai sản phẩm thiết kế có thể triển khai trực tiếp là cơ sở dữ liệu và thiết kế chaincode.

Về phía \textbf{đặt bài toán và phân tích yêu cầu} (mục 2.1--2.2), chương đã xác lập ba mục tiêu cốt lõi -- \textit{bỏ phiếu duy nhất một lần}, \textit{bí mật lá phiếu} và \textit{tính kiểm chứng của kết quả} -- cùng sáu tính chất hệ thống bắt buộc: bí mật, duy nhất, xác thực, không thể chối bỏ, kiểm chứng và ẩn danh. Đây không phải sáu tính chất độc lập mà là một tập ràng buộc phụ thuộc lẫn nhau; việc thiếu bất kỳ tính chất nào cũng làm suy giảm nghiêm trọng nền tảng tin cậy của cuộc bầu cử. Ba tác nhân chính -- quản trị viên, cử tri và khách -- được xác định rõ vai trò và quyền hạn, tạo cơ sở để thiết kế phân quyền chi tiết ở các tầng tiếp theo.

Về phía \textbf{mô hình hóa ca sử dụng và biểu đồ tuần tự} (mục 2.3--2.5), chương đã triển khai từ biểu đồ ca sử dụng tổng quát đến các đặc tả chi tiết cho từng ca sử dụng, bao gồm tiền điều kiện, luồng chính, luồng ngoại lệ và hậu điều kiện. Các biểu đồ tuần tự phác thảo rõ mối quan hệ giao tiếp giữa các thành phần trong ba giai đoạn then chốt: xác thực và lấy chữ ký mù tập thể, nộp phiếu lên blockchain, và kiểm chứng phiếu sau bầu cử.

Về phía \textbf{thiết kế dữ liệu} (mục 2.6--2.7), hai sản phẩm thiết kế được đưa ra tương ứng với hai tầng lưu trữ của hệ thống:

\begin{itemize}
    \item \textbf{Cơ sở dữ liệu off-chain} mô hình hóa toàn bộ trạng thái vòng đời cuộc bầu cử -- từ thông tin tổ chức, cử tri, ứng viên đến bản ghi phiếu mù và biên lai xác minh -- với các ràng buộc cứng ngăn chặn trùng lặp và đảm bảo tính nhất quán trước khi bất kỳ dữ liệu nào được đẩy lên blockchain.

    \item \textbf{Chaincode \textit{VoteLedgerContract}} thiết kế theo nguyên tắc tối giản: không quản lý vòng đời hay xác thực danh tính, chỉ đóng vai trò tầng \textit{chứng cứ bất biến}. Ba nhóm hợp đồng -- \textit{Vote Contract}, \textit{Merkle Contract} và \textit{Reveal Contract} -- tương ứng với ba pha bắt buộc trên blockchain: nhận phiếu, cố định gốc Merkle khi đóng bầu cử, và giải mù để kiểm phiếu. Mười hàm chaincode được thiết kế rõ ràng về bất biến kiểm tra, điều kiện tiên quyết và kết quả ghi nhận, đảm bảo tính không thể chối bỏ và khả năng kiểm toán độc lập.
\end{itemize}

Điểm xuyên suốt toàn bộ chương là nguyên tắc \textbf{bảo vệ đa tầng}: không một tầng đơn lẻ nào -- dù là cơ sở dữ liệu hay chaincode -- phải tự mình gánh chịu toàn bộ trách nhiệm bảo mật. Thay vào đó, mỗi tầng áp đặt một tập ràng buộc riêng, và sự kết hợp của các tầng mới tạo nên đảm bảo tổng thể cho sáu tính chất hệ thống đã đề ra.

Các kết quả thiết kế trong chương 2 -- cụ thể là cơ sở dữ liệu , giao thức luồng nghiệp vụ và đặc tả chaincode -- sẽ được hiện thực hóa trực tiếp thành mã nguồn và cấu hình triển khai trong chương 3.
